% !TEX encoding = UTF-8 Unicode
% !TEX root = thesis-main.tex


\chapterr{Λειτουργίες του \textlatin{MeL}}
Στο προηγούμενο κεφάλαιο παρουσιάσαμε αναλυτικά τη Βάση Δεδομένων στην οποία στηρίζεται το \textlatin{MeL} καθώς επίσης και τις απαιτήσεις για την εγκατάσταση και τη λειτουργία του σε τοπικό επίπεδο. Σε αυτό το κεφάλαιο θα παρουσιάσουμε τις λειτουργίες της πλατφόρμας. Πιο συγκεκριμένα θα παραθέσουμε αναλυτικά εικόνες από την πλατφόρμα και επεξηγήσεις για κάθε σενάριο λειτουργίας της.\section{Πρώτη επαφή}
Σε αυτήν την ενότητα θα δούμε για πρώτη φορά την πλατφόρμα και θα αναλύσουμε τις επιλογές των χρηστών-επισκεπτών για να μπορέσουν να διεισδύσουν περισσότερο μέσα σε αυτή.

\subsection{Αρχική σελίδα}
Στο Σχήμα \ref{fig:start_page} βλέπουμε την αρχική σελίδα σαν χρήστες-επισκέπτες, όπου μπορούμε να δούμε συνοπτικά το θέμα με το οποίο ασχολείται η πλατφόρμα και τις βασικές της λειτουργίες. Το θέμα αρχικά μπορεί να εξαχθεί από τον τίτλο της πλατφόρμας \textlatin{MeL (Medical e-Learning),} ηλεκτρονική εκπαίδευση ιατρικού/φαρμακευτικού ενδιαφέροντος. Διακρίνουμε Άρθρα και Τεστ στα οποία δεν έχουμε πρόσβαση ακόμα. Δεξιά υπάρχει μία στήλη όπου μπορούμε να κάνουμε είσοδο στη πλατφόρμα αν έχουμε δημιουργήσει λογαριασμό, αλλιώς μπορούμε να δημιουργήσουμε έναν κάνοντας κλικ στο κουμπί Εγγραφή.

\begin{figure}[tbh]
\centering
% \includegraphics[scale=0.45]{figures/ch3-start_page}
\caption{Αρχική σελίδα \textlatin{MeL.}}
\label{fig:start_page}
\end{figure}

\subsection{Δημιουργία λογαριασμού}
Έχοντας κάνει κλικ στο κουμπί Εγγραφή στην αρχική σελίδα θα προωθηθούμε στη σελίδα δημιουργίας λογαριασμού (Φόρμα Εγγραφής), η οποία είναι υπεύθυνη για την δημιουργία λογαριασμού στη πλατφόρμα. Συμπληρώνουμε υποχρεωτικά όλα τα πεδία ξεκινώντας από τα στοιχεία του λογαριασμού του χρήστη. Το Όνομα Χρήστη θα είναι το όνομα το οποίο θα δίνει ο χρήστης για να εισέλθει στη πλατφόρμα. Επίσης, θα δώσει και τον κωδικό τον οποίο θα χρησιμοποιεί για να εισέρχεται στο πεδίο Κωδικός Πρόσβασης. Θα ξανά πληκτρολογήσει τον κωδικό αυτό στο πεδίο Επαλήθευση Κωδικού Πρόσβασης για την αποφυγή πιθανού λάθους στην πληκτρολόγηση αυτού. Τελευταίο πεδίο για τα στοιχεία του λογαριασμού του χρήστη είναι το \textlatin{email.} Τέλος, ζητούνται από το χρήστη δύο προσωπικά στοιχεία για την ολοκλήρωση της εγγραφής του, το όνομά του και το επίθετό του στα αντίστοιχα πεδία (βλ. Σχήμα \ref{fig:register}). Πατώντας το κουμπί Υποβολή, ελέγχεται η ορθότητα των στοιχείων που έδωσε ο χρήστης, και αν όλα είναι σωστά ολοκληρώνεται η φόρμα εγγραφής και αποστέλλεται μήνυμα επιβεβαίωσης στο \textlatin{email} του χρήστη, το οποίο καλείται ο χρήστης να ανοίξει και να ακολουθήσει το σύνδεσμο που του έχει αποσταλεί ώστε να ενεργοποιήσει το λογαριασμό του. Σε περίπτωση που κάποιο από τα πεδία είναι λάθος, εμφανίζεται κατάλληλο μήνυμα για την επισήμανση αυτού. Με την επιτυχή ολοκλήρωση εγγραφής ο χρήστης-επισκέπτης θα μπορέσει να κάνει είσοδο στη πλατφόρμα ως χρήστης-μέλος.

\begin{figure}[tbh]
\centering
% \includegraphics[scale=0.50]{figures/ch3-register}
\caption{Δημιουργία λογαριασμού.}
\label{fig:register}
\end{figure}

\section{Λειτουργίες μελών της πλατφόρμας}

\subsection{Είσοδος στην πλατφόρμα}

\subsection{Αρχική σελίδα μέλους}
Σε αυτή την ενότητα βλέπουμε την επιτυχή είσοδο του χρήστη-μέλους και τις καινούργιες λειτουργίες του συστήματος που εμφανίστηκαν, όπως οι επιλογές για τη ρύθμιση του προφίλ του χρήστη, τους φίλους, τα μηνύματα, τις ειδοποιήσεις και τα αιτήματα φιλίας λειτουργίες για τις οποίες θα μιλήσουμε στις επόμενες ενότητες του κεφαλαίου (βλ. Σχήμα \ref{fig:start_page_user}). 

\begin{figure}[tbh]
% \centerline{\includegraphics[scale=0.45]{figures/ch3-start_page_user}}
\caption{Αρχική σελίδα μέλους.}
\label{fig:start_page_user}
\end{figure}

Ο χρήστης μπορεί να περιηγηθεί στη πλατφόρμα ανοίγοντας ένα Άρθρο ή ένα Τεστ ή τις αντίστοιχες καρτέλες. Επίσης, μπορεί να τα αναζητήσει κιόλας από το πεδίο αναζήτηση καθώς και άλλους χρήστες με τις αντίστοιχες επιλογές του στην περιοχή της αναζήτησης (βλ. Σχήματα \ref{fig:search_test} και \ref{fig:search_user}).


\begin{figure}[tbh]
	\begin{minipage}[t]{0.45\linewidth}
	\centering
    	% \includegraphics[scale=0.60]{figures/ch3-search_test}
	\caption{Αναζήτηση τεστ με τον όρο \textquotedblleft καρδ\textquotedblright.}
	\label{fig:search_test}	
	\end{minipage}\hfill
	\begin{minipage}[t]{0.45\linewidth}
	\centering
	% \includegraphics[scale=0.7]{figures/ch3-search_user.png}
	\caption{Αναζήτηση χρήστη με τον όρο \textlatin{\textquotedblleft gian\textquotedblright .}}
	\label{fig:search_user}
	\end{minipage} 
\end{figure}


\subsection{Επεξεργασία προφίλ}


\subsection{...}

\section{Λειτουργίες διαχειριστών της πλατφόρμας}


\subsection{Δημιουργία/Επεξεργασία κατηγορίας}

\subsection{...}	



\section{Λειτουργίες υπέρ-διαχειριστή της πλατφόρμας}



Σε αυτό το κεφάλαιο είδαμε τις περιπτώσεις χρήσεις της πλατφόρμας από όλους τους χρήστες για όλες τις λειτουργίες με ζωντανά παραδείγματα. Όλα αυτά όμως για το μέρος της εργασίας που αφορά την ιστοσελίδα. Στο επόμενο κεφάλαιο θα δούμε το αντίστοιχο μέρος της για την λειτουργία της στο λειτουργικό σύστημα \textlatin{Android.}